\section{Grundlagen}
\label{sec:grundlagen}

Um das Thema tiefgreifend analysieren zu können, ist es notwendig, die mathematischen und physikalischen Prinzipien zu verstehen.

\section{Mathematische Formeln}

Die Navier-Stokes-Gleichungen sind ein fundamentales Konzept der Strömungsmechanik. Gleichung \ref{eq:navier_stokes} zeigt die inkompressiblen Navier-Stokes-Gleichungen in Tensornotation.

\begin{equation}
\label{eq:navier_stokes}
\frac{\partial u_i}{\partial t} + u_j \frac{\partial u_i}{\partial x_j} = - \frac{1}{\rho} \frac{\partial P}{\partial x_i} + \nu \frac{\partial^2 u_i}{\partial x_j^2}
\end{equation}

Solche komplexen Formeln lassen sich leicht darstellen und referenzieren. Hierbei wird ein Ansatz verwendet, der sehr stark an die Schreibweise auf dem Papier erinnert, aber intern wie ein modernes Programmierframework funktioniert.

Ghia et al. haben diese Gleichungen intensiv untersucht und Lösungsansätze entwickelt, um Strömungen bei hohen Reynolds-Zahlen zu berechnen \cite{ghia_1982}.
